\section{Inner Monologue：从 Open Loop 到 Feedback-Conditioned Replanning}

SayCan 每一步会重新选择技能，却可能不知道上一步是否成功。若机器人没抓到可乐仍继续“拿给用户”，高层计划在语言上连贯、在世界中却已经失效。Inner Monologue 为 planner 增加 scene description、clarification 和 success detector，让语言推理真正闭环。

\subsection{Agenda 与 Open-Loop Failure}

Agenda 进入 Inner Monologue。下一张 slide 用失败轨迹说明问题：用户要一杯饮料，planner 选择可乐；抓取失败后，系统仍按原序列继续移动并宣告完成。缺少 observation feedback 时，CoT 只是 open-loop script。

\lecturefigure{slide-30-agenda-inner-monologue.jpg}{Agenda：从 SayCan 转向带环境反馈的 Inner Monologue}{官方录像恢复幻灯片 V030；Stanford CS25 Lecture 15}

\lecturefigure{slide-31-saycan-open-loop-failure.jpg}{Open-loop failure：抓取可乐失败后，planner 仍继续执行旧计划}{官方录像恢复幻灯片 V031；Huang et al., 2022}

\begin{importantbox}{具身 Planning 必须维护可更新状态}
语言计划隐含一个 belief：桌上有什么、机器人拿到了什么、用户偏好是什么、当前子目标是否完成。任何 action 后都可能改变 belief。若系统不读取新观测，后续 token 再合理也只是对过期世界状态的推理。
\end{importantbox}

\subsection{Passive Scene Description}

上一小节的失败轨迹表明，open-loop planner 即使生成了合理步骤，也会在一次抓取失败后继续沿过期计划前进。最直接的修复不是立刻改变 planner，而是先让环境状态持续回到语言上下文。本节考察 passive feedback 如何把检测结果和任务进展压缩成文本；看图时要同时关注它提供了哪些状态，以及哪些物理事实因 detector ontology 不足而被丢弃。

被动反馈由 detector 或 VLM 持续写入，例如“我看到一罐可乐和一瓶青柠汽水”“黄色积木已在黄色碗中”。这些描述把高维图像压成 planner 可消费的文本状态，使 LLM 在每步后知道世界是否发生预期变化。

\lecturefigure{slide-32-inner-monologue-passive-feedback.jpg}{Passive feedback：object detection、task-progress description 与 visual QA}{官方录像恢复幻灯片 V032；Huang et al., 2022}

\begin{knowledgebox}{Passive feedback 的信息瓶颈}
Scene caption 只保留 detector 会表达的概念。若系统没有“杯子倾倒”“液体泄漏”或“夹爪空了”的标签，planner 看不到这些事实。文本化降低接口复杂度，也把 perception ontology 变成上限。
\end{knowledgebox}

\subsection{Active Scene Description 与 Clarification}

主动反馈只在 planner 不确定时请求，例如询问用户偏好、澄清指令或查询物体位置。slide 33 把 active scene description 与 visual QA in context 放在一起：机器人可以问“你想要哪种饮料”，再根据回答更新计划，而不是在不确定条件下猜测。

\lecturefigure{slide-33-inner-monologue-active-feedback.jpg}{Active feedback：只在需要时询问用户、环境或视觉问答模块}{官方录像恢复幻灯片 V033；Huang et al., 2022}

若把文本 belief 记为 $b_t$，执行观测为 $o_{t+1}$，success signal 为 $z_{t+1}$，用户反馈为 $u_{t+1}$，闭环更新可抽象为
\begin{equation}
b_{t+1}=U(b_t,o_{t+1},z_{t+1},u_{t+1}),
\qquad
k_{t+1}=\arg\max_k \operatorname{Score}(k\mid g,b_{t+1}).
\end{equation}
与 open-loop plan 不同，每一步选择都基于更新后的 belief。

\subsection{Success Detector 与结果表}

结果 slide 上半部比较 CLIPort、oracle、LLM、object feedback 与 Inner Monologue 等组合，下半部给出具体对话和 task-progress detector。读表时应先区分 seen / unseen tasks，再比较只加 object feedback、只加 success feedback 与二者结合的差异。

\lecturefigure{slide-34-inner-monologue-results.jpg}{Inner Monologue results：object feedback、success detector 与闭环 replanning}{官方录像恢复幻灯片 V034；Huang et al., 2022}

\begin{warningbox}{Feedback 不是自动真实}
Success detector 可能误报，scene description 可能漏物体，用户回答也可能含糊。Closed loop 只是让错误有机会被纠正；若反馈模型系统性偏差，planner 会在错误 belief 上反复自洽。
\end{warningbox}

\subsection{Importing Common Sense 的边界}

最后一张 slide 用 2018 年“机器人瓶颈是 high-level semantic planning”的说法，与 2022 年 LLM 的“no”形成对比。下方 lesson 更准确：若 language 是系统 universal API，foundation models 可以注入 common sense。它降低高层语义规划成本，却没有消除 perception、skill 和 control bottlenecks。

\lecturefigure{slide-35-foundation-common-sense-lesson.jpg}{利用 foundation models 注入 common sense，但不替代物理系统}{官方录像恢复幻灯片 V035；Stanford CS25 Lecture 15}

\teachervoice{讲者没有把 LLM 说成机器人“大脑”。他的论点是 language interface 使外部模型的 common sense 能进入已有系统；如果 detector、affordance 或 controller 不提供正确接口，LLM 再强也没有可操作的世界状态。}

\begin{lstlisting}[caption={SayCan + Inner Monologue 的 closed-loop planner},label={lst:closed-loop-planner}]
belief = describe_scene(initial_observation)
history = []

while not task_complete(belief):
    candidates = language_model.propose_skills(goal, belief, history)
    scored = []
    for skill in candidates:
        useful = language_model.score(skill, goal, history)
        possible = affordance_model.score(skill, belief)
        scored.append((useful * possible, skill))

    skill = max(scored)[1]
    outcome = execute(skill)
    feedback = collect_scene_success_and_user_feedback(outcome)
    belief = update_belief(belief, feedback)
    history.append((skill, feedback))
\end{lstlisting}

\subsection{本章小结}

Inner Monologue 把语言 planning 从 open-loop sequence 变成 observation-conditioned loop。Passive scene、active clarification 和 success detection 提供不同反馈通道；系统能力仍受 perception ontology 和 detector reliability 限制，但至少能在世界偏离计划时重新规划。

